ابتدا واکنش کلی پیل و پتانسیل استاندارد آن را مشخص می‌کنیم. سپس با استفاده از معادله نرنست (\lr{Nernst Equation}) و جایگذاری مقادیر داده شده، غلظت مجهول یون کادمیوم را محاسبه می‌کنیم.

\begin{flalign*}
    & \text{واکنش آند (اکسایش):} \quad Zn \rightarrow Zn^{2+} + 2e^- \qquad E^\circ_{ox} = +0.76 \text{ V} &\\
    & \text{واکنش کاتد (کاهش):} \quad Cd^{2+} + 2e^- \rightarrow Cd \qquad E^\circ_{red} = -0.40 \text{ V} &\\
    & \text{واکنش کلی:} \quad Zn + Cd^{2+} \rightleftharpoons Zn^{2+} + Cd \qquad n = 2 &\\
    & E^\circ_{\text{cell}} = E^\circ_{ox} + E^\circ_{red} = 0.76 + (-0.40) = 0.36 \text{ V} &\\
    & E_{\text{cell}} = E^\circ_{\text{cell}} - \frac{0.0592}{n} \log \frac{[Zn^{2+}]}{[Cd^{2+}]} &\\
    & 0.40 = 0.36 - \frac{0.0592}{2} \log \frac{0.09}{[Cd^{2+}]} &\\
    & 0.04 = -0.0296 \log \frac{0.09}{[Cd^{2+}]} &\\
    & \log \frac{0.09}{[Cd^{2+}]} = \frac{0.04}{-0.0296} \approx -1.351 &\\
    & \frac{0.09}{[Cd^{2+}]} = 10^{-1.351} \approx 0.0445 &\\
    & [Cd^{2+}] = \frac{0.09}{0.0445} \approx 2.02 \text{ M} &
\end{flalign*}